\subsection{Front-end}
\label{subsec:architettura_frontend}
\glossario{Angular} adotta il pattern architetturale Model-View-ViewModel (MVVM), un pattern architetturale che separa le responsabilità dell'interfaccia utente (UI) dalla \textit{business logic} di un'applicazione.
In \glossario{Angular}, l'MVVM è implementato attraverso i componenti, che fungono da ViewModel e gestiscono il legame tra il Model (dati e logica dell'applicazione) e la View (interfaccia utente).
Grazie al data binding bidirezionale, \glossario{Angular} sincronizza automaticamente i cambiamenti tra il Model e la View, semplificando l'interazione tra i livelli dell'applicazione.
Il pattern architetturale MVVM è formato dai seguenti tre layer:
\begin{itemize}
    \item \textbf{Model}: Si riferisce ai dati dell'applicazione e alla \textit{business logic}, come la validazione dei dati o l'elaborazione delle operazioni. 
    Il Model non è a conoscenza dei livelli View o ViewModel, il che lo rende indipendente dalla presentazione. 
    In un'architettura MVVM, il Model rappresenta generalmente i dati che l'applicazione elabora e visualizza.
    In \glossario{Angular}, il Model è generalmente costituito da classi TypeScript che definiscono la struttura e il comportamento dei dati dell'applicazione. 
    Queste classi possono includere proprietà, metodi e regole di validazione;
    \item \textbf{View}: È responsabile della struttura, del layout e dell'aspetto dell'interfaccia utente. 
    La View riceve le interazioni dell'utente e le passa al ViewModel tramite data binding. 
    Non contiene la \textit{business logic}, ma si collega al ViewModel, che fornisce i dati e le operazioni necessarie per visualizzare e interagire con le informazioni.
    La View in \glossario{Angular} è rappresentata dai template di \glossario{Angular}, che sono file HTML arricchiti con direttive e binding, questi template definiscono come i dati vengono presentati all'utente e come vengono catturati gli input dell'utente.
    Quando i dati del Model cambiano a seguito dell'elaborazione della \textit{business logic}, la View si aggiorna automaticamente grazie al data binding bidirezionale con il ViewModel;
    \item \textbf{ViewModel}: Il ViewModel funge da livello intermedio tra la View e il Model.
    Il ViewModel rappresenta lo stato dei dati presenti nel Model e gestisce la comunicazione bidirezionale tra il Model e la View attraverso la tecnica del data binding.
    Inoltre, il ViewModel può contenere logica legata all'elaborazione e alla validazione dei dati, consentendo operazioni più complesse, come il filtraggio, l'ordinamento o il raggruppamento dei dati, senza dover intervenire sul livello della View o del Model.
    In \glossario{Angular}, il ViewModel è spesso rappresentato dai Component e dalle relative classi TypeScript. 
    In \glossario{Angular} i Component contengono l'\textit{application logic}, espongono proprietà e metodi utilizzati nei template per il data binding e comunicano con il Model attraverso i servizi (indicati con il decoratore \textit{@Injectable}). 
\end{itemize}

I concetti fondamentali implementati da \glossario{Angular} sono:
\begin{itemize}
    \item \textbf{Data Binding}: \glossario{Angular} utilizza un avanzato sistema di data binding che consente una sincronizzazione automatica tra il Model (dati dell'applicazione) e la View (interfaccia utente).
    Il data binding può avvenire in una o due direzioni, con il two-way data binding i dati rimangono sincronizzati tra Model e View, garantendo un aggiornamento dinamico dell'interfaccia utente;
    \item \textbf{Modularità}:  \glossario{Angular} è basato su un'architettura modulare, in cui un'applicazione è suddivisa in moduli (\textit{NgModules}).
    Ogni modulo può contenere componenti, servizi e altre funzionalità, permettendo una suddivisione chiara e organizzata del codice, facilitando lo sviluppo e il riutilizzo delle componenti;
    \item \textbf{Component-based architecture}:  \glossario{Angular} segue un'architettura basata sui componenti, in cui l'interfaccia utente è costruita attraverso componenti riutilizzabili e indipendenti.
    Ogni componente gestisce la propria logica e il proprio template, migliorando la manutenibilità e la scalabilità dell'applicazione;
    \item \textbf{Directive e template engine}: \glossario{Angular} estende l'HTML con directive (direttive) personalizzate che permettono di manipolare il DOM in modo dinamico.
    Esistono directive strutturali per il controllo del flusso della UI e directive attributo per la modifica degli elementi HTML;
    \item \textbf{RxJS e programmazione reattiva}: \glossario{Angular} utilizza \textit{RxJS} (Reactive Extensions for \glossario{Javascript}) per gestire lo stato e la comunicazione tra componenti in modo asincrono e reattivo.
    Questo permette di gestire flussi di dati come eventi utente, richieste HTTP e aggiornamenti dello stato attraverso Observable e operatori reattivi.
\end{itemize}
Il codice sviluppato per il front-end verrà organizzato nel seguente modo:
\iffalse
    frontend/
    │── src/
        ├── app/
        │   ├── core/                     # Moduli e servizi condivisi (singleton)
        │   │   ├── services/             # Servizi globali (es. autenticazione, API)
        │   │   │   ├── auth.service.ts
        │   │   │   ├── http.service.ts
        │   │   ├── interceptors/         # Interceptor per HTTP (es. JWT, logging)
        │   │   │   ├── auth.interceptor.ts
        │   │   ├── guards/               # Route guards (es. autenticazione)
        │   │   │   ├── auth.guard.ts
        │   │   ├── core.module.ts        # Modulo Core (singleton)
        │   ├── shared/                   # Moduli, componenti e pipe riutilizzabili
        │   │   ├── components/           # Componenti condivisi (es. pulsanti, modali)
        │   │   │   ├── navbar/           # Esempio di un componente riutilizzabile
        │   │   │   │   ├── navbar.component.ts
        │   │   │   │   ├── navbar.component.html
        │   │   ├── directives/           # Direttive personalizzate (es. autofocus)
        │   │   ├── pipes/                # Pipe personalizzate (es. formattazione date)
        │   │   ├── shared.module.ts      # Modulo Shared (componenti riutilizzabili)
        │   ├── features/                 # Moduli feature-based
        │   │   ├── users/                # Feature: gestione utenti
        │   │   │   ├── models/           # Modelli dati (User.ts, Role.ts)
        │   │   │   ├── services/         # Servizi specifici della feature
        │   │   │   ├── components/       # Componenti della feature
        │   │   │   ├── pages/            # Pagine della feature (es. lista utenti)
        │   │   │   ├── users.module.ts   # Modulo utenti
        │   ├── layout/                   # Layout generali (header, sidebar)
        │   ├── app.component.ts          # Root component
        │   ├── app.module.ts             # Root module
        ├── assets/                       # File statici (es. immagini, JSON)
        ├── environments/                 # Configurazioni per env (dev, prod)
        ├── main.ts                        # Bootstrap Angular
        ├── styles.scss                    # Stili globali
\fi
\dirtree{%
.1 ArtificialQI/.
.2 src/.
.3 app/.
.4 core/.
.5 services/.
.5 models/.
.5 interceptors/.
.4 shared/.
.5 components/.
.5 directives/.
.5 pipes/.
.4 features/.
.5 feature1/.
.6 components/.
.6 pages/.
.4 layout/.
.2 assets/.
.2 environments/.
}
Di seguito è riportata una breve descrizione delle cartelle principali:
\begin{itemize}
    \item \textbf{src}: contiene il codice sorgente;
    \item \textbf{assets}: contiene i file statici, come immagini o font;
    \item \textbf{environments}: contiene file di configurazione per ambienti diversi (sviluppo e produzione).
\end{itemize}
La cartella \textit{src} contiene il codice sorgente dell'applicazione all'interno della sottocartella \textit{app}, \textit{app} è suddivisa in:
\begin{itemize}
    \item \textbf{core}: contiene modelli e \textit{business logic} all'interno delle sottocartelle:
    \begin{itemize}
        \item \textbf{models}: interfacce e classi che rappresentano i dati e le entità principali dell'applicazione;
        \item \textbf{services}: gestiscono la \textit{business logic} e le chiamate API;
        \item \textbf{interceptors}: modificano le chiamate HTTP (ad esempio per aggiungere autenticazione o logging) e intercettano errori API per gestirli in modo centralizzato.
    \end{itemize}
    \item \textbf{shared}: contiene moduli, componenti e pipe riutilizzabili, all'interno delle sottocartelle:
    \begin{itemize}
        \item \textbf{components}: componenti condivise (es. pulsanti);
        \item \textbf{directives}: direttive personalizzate;
        \item \textbf{pipes}: pipe personalizzate per formattazione dati.
    \end{itemize} 
    \item \textbf{features}: moduli feature-base, ogni feature ha una cartella dedicata organizzata in:
    \begin{itemize}
        \item \textbf{components}: contiene i componenti specifici della feature, che possono essere combinati per formare una pagina;
        \item \textbf{pages}: contiene le pagine principali della feature, ogni pagina funge da container component del ViewModel nel pattern MVVM.
    \end{itemize}
    \item \textbf{layout}: contiene componenti per il layout dell'app (es. navbar, sidebar, footer).
\end{itemize}
